\documentclass[11pt]{article}
\usepackage[margin=1in]{geometry}
\usepackage{amsmath,amssymb}
\usepackage{graphicx}
\usepackage{booktabs}

\title{Burn-in and finite-time flux-tail diagnostic}
\author{}
\date{}

\begin{document}
\maketitle

\section*{1 \quad Burn-in time}

\subsection*{Setup}

\begin{itemize}
    \item Closed chain, $T_1=10$, $T_n=2$, $\gamma=0.1$.
    \item Conserved action current across the middle bond
    \[
        J_m = 4I_{m-1}I_m
        \sin\!\big(2(\varphi_m-\varphi_{m-1})\big).
    \]
    \item Time step $dt=5\times 10^{-4}$.
    \item No burn-in.  Start from the deterministic initial condition used in
    the canonical current code.
    \item For $n=10,20,40,80$, simulate $1024$ trajectories and record
    profiles and flux at times $t=50,100,150,\ldots,500$.
\end{itemize}

\subsection*{Results}

\begin{table}[h]
\centering
\begin{tabular}{rrrrr}
\toprule
$n$ & stable profile time & stable flux time & final cumulative $J$ & last-window $J$ at $t=500$ \\
\midrule
10 & $\sim 100$ & $\sim 100$ & 0.3958 & 0.3978 \\
20 & $100$--$200$ & $100$--$200$ & 0.1175 & 0.1230 \\
40 & $200$--$300$ & $200$--$300$ & 0.0282 & 0.0302 \\
80 & $\gtrsim 200$ & not fully resolved & 0.0057 & 0.0061 \\
\bottomrule
\end{tabular}
\end{table}

\begin{figure}[h]
\centering
\includegraphics[width=\textwidth]{analysis/burnin_terminal_profile.pdf}
\caption{No-burn-in terminal action profiles.  The curves show
$\mathbb{E}[I_j(t)]$ at $t=50,100,\ldots,500$.}
\end{figure}

\begin{figure}[h]
\centering
\includegraphics[width=\textwidth]{analysis/burnin_flux_timeseries.pdf}
\caption{No-burn-in current relaxation.  Left: cumulative average over
$[0,t]$.  Right: average over the last interval of length $50$.}
\end{figure}

\begin{itemize}
    \item The relaxation time grows with chain length.
    \item For $n=10,20$, both profile and current stabilize quickly.
    \item For $n=40$, the current is still biased low at early times; by
    about $t=200$--$300$ the last-window current is close to the later values.
    \item For $n=80$, $t=50$ is clearly too short.  The profile shape becomes
    much closer to later profiles after about $t=150$--$200$, but the current
    is small and noisy.
    \item For the flux-tail pilot below, we therefore use the conservative
    common burn-in $B=500$ for $n\le 40$.
\end{itemize}

\section*{2 \quad Probability distribution of flux}

\subsection*{Setup}

\begin{itemize}
    \item Same closed-chain parameters: $T_1=10$, $T_n=2$, $\gamma=0.1$.
    \item Burn-in $B=500$ from the diagnostic above.
    \item For $n=10,20,30,40$, simulate $1024$ trajectories.
    \item Measurement length $200$, divided into ten blocks of length $20$.
    \item Construct prefix averages
    \[
        \overline J_\tau
        = \frac{1}{\tau}\int_B^{B+\tau} J_m(t)\,dt,
        \qquad
        \tau=20,40,\ldots,200.
    \]
    \item For thresholds $A=0.01,0.02,\ldots$, sort the samples and estimate
    $P[\overline J_\tau>A]$.
\end{itemize}

\subsection*{Results at $\tau=200$}

\begin{table}[h]
\centering
\begin{tabular}{rrrrrc}
\toprule
$n$ & mean $J$ & std & $\lambda$ & $R^2$ & fit window in $A$ \\
\midrule
10 & 0.3959 & 0.0607 & 31.4  & 0.9918 & $[0.45,\,0.55]$ \\
20 & 0.1175 & 0.0303 & 62.0  & 0.9909 & $[0.15,\,0.19]$ \\
30 & 0.0539 & 0.0204 & 91.0  & 0.9955 & $[0.08,\,0.10]$ \\
40 & 0.0298 & 0.0163 & 126.1 & 1.0000 & $[0.05,\,0.07]$ \\
\bottomrule
\end{tabular}
\end{table}

Here $\lambda$ is the descriptive slope magnitude from a linear fit
\[
    \log P[\overline J_\tau>A] \approx -\lambda A+c
\]
over the listed upper-middle-tail window.  It is not a transport exponent.

\begin{figure}[h]
\centering
\includegraphics[width=\textwidth]{analysis/tau_tail_survival_n40.pdf}
\caption{Example for $n=40$: survival plots
$\log P[\overline J_\tau>A]$ versus threshold $A$ for
$\tau=20,40,\ldots,200$.}
\end{figure}

\begin{figure}[h]
\centering
\includegraphics[width=0.72\textwidth]{analysis/tau_logprob_vs_tau_n40.pdf}
\caption{Example for $n=40$: fixed-threshold dependence on $\tau$.}
\end{figure}

\begin{itemize}
    \item For fixed $\tau$, $\log P[\overline J_\tau>A]$ decreases with
    threshold $A$.
    \item The upper-middle tail is locally close to linear in $A$ on the log
    survival plot, but the full distribution is not exponential.
    \item As $\tau$ increases, the time-averaged current distribution narrows.
    \item If $A$ is below the mean current, $P[\overline J_\tau>A]$ can
    increase with $\tau$ because the distribution concentrates near a positive
    mean.
    \item If $A$ is above the mean current, $P[\overline J_\tau>A]$ decreases
    with $\tau$.
    \item The rate proxy
    \[
        -\frac{1}{\tau}\log P[\overline J_\tau>A]
    \]
    does not yet show a clean collapse across $\tau$ in this $1024$-trajectory
    pilot.  Therefore these plots should be described as finite-time
    flux-tail diagnostics, not as a large-deviation result.
\end{itemize}

\section*{Takeaway}

\begin{itemize}
    \item Burn-in matters, especially for longer chains.
    \item The diagnostic suggests that $B=500$ is a reasonable conservative
    burn-in for the present $n\le 40$ finite-time tail pilot.
    \item The finite-time flux tails have locally exponential-looking
    upper-middle regions, but a large-deviation claim would require many more
    trajectories and clearer $\tau$-scaling.
    \item A $10^5$-trajectory run is the natural next step if we want to
    resolve smaller probabilities reliably.
\end{itemize}

\end{document}
